\subsection{Atomicita změn a code review}
\label{subsec:kelvin-development-atomicity}

Jedním z hlavních pravidel vývoje, zdokumentovaných v příručce pro rozšiřování systému\footnote[2]{\href{https://mrlvsb.github.io/kelvin/developers-guide/development\#contributing}{https://mrlvsb.github.io/kelvin/developers-guide/development\#contributing}}, je požadavek na malé a atomické pull requesty.
Každý pull request by měl obsahovat maximálně $\sim500$ změněných řádků kódu a měl by se zaměřovat pouze na jednu konkrétní změnu nebo problém.

Cílem tohoto přístupu je zjednodušení kontrolního procesu (\textit{code review}) a snížení rizika zavlečení chyb do produkčního systému.
V praxi to však znamená nutnost rozdělovat i logicky související úpravy do více menších kroků, což může celý vývojový proces prodlužovat.

\endinput
